% 八点周期模糊核(0.2,0.6,0.2)，H(w)=0.6+0.4cos(w)。三条曲线为逐频率增益的解析值：
% 逆滤波1/H；常数比维纳 H/(H^2+0.05)；一阶差分约束 H/(H^2+0.05(2-2cos w))。宽约95mm。
\begin{tikzpicture}[x=72mm,y=8.5mm,font=\figurefont\small]
 \draw[BookMuted,line width=.7pt,->] (0,0)--(1.08,0) node[right] {$\omega$};
 \draw[BookMuted,line width=.7pt,->] (0,0)--(0,5.5) node[above] {增益};
 \foreach \y in {1,2,3,4,5} {
  \draw[BookRule,line width=.5pt] (0,\y)--(1,\y);
  \node[anchor=east] at (-.01,\y) {$\y$};
 }
 \node[anchor=east] at (-.01,0) {$0$};
 \foreach \x/\lab in {0/0,.5/{\pi/2},1/\pi} {
  \draw[BookMuted,line width=.5pt] (\x,0)--(\x,-.1);
  \node[below] at (\x,-.1) {$\lab$};
 }
 \draw[BookAmber,dashed,line width=1pt] plot coordinates {
  (0,1)(.0833,1.014)(.1667,1.057)(.25,1.133)(.3333,1.25)(.4167,1.421)
  (.5,1.667)(.5833,2.014)(.6667,2.5)(.75,3.153)(.8333,3.943)(.9167,4.681)(1,5)};
 \draw[BookBlue,line width=1pt] plot coordinates {
  (0,.952)(.0833,.964)(.1667,1.001)(.25,1.064)(.3333,1.159)(.4167,1.291)
  (.5,1.463)(.5833,1.675)(.6667,1.905)(.75,2.106)(.8333,2.218)(.9167,2.234)(1,2.222)};
 \draw[BookTeal,line width=1pt] plot[mark=*,mark size=1.2pt] coordinates {
  (0,1)(.0833,1.01)(.1667,1.041)(.25,1.092)(.3333,1.159)(.4167,1.236)
  (.5,1.304)(.5833,1.333)(.6667,1.29)(.75,1.169)(.8333,1.011)(.9167,.882)(1,.833)};
 \draw[BookAmber,dashed,line width=1pt] (.04,6.2)--(.12,6.2);
 \node[anchor=west] at (.13,6.2) {逆滤波};
 \draw[BookBlue,line width=1pt] (.34,6.2)--(.42,6.2);
 \node[anchor=west] at (.43,6.2) {常数比维纳};
 \draw[BookTeal,line width=1pt] (.71,6.2)--(.79,6.2);
 \node[anchor=west] at (.8,6.2) {差分约束};
\end{tikzpicture}
